\documentclass{article}
\usepackage[UTF8]{ctex}
\usepackage{amsmath,amssymb,amsfonts}
\usepackage{enumitem}
\usepackage{geometry}
\geometry{a4paper, margin=1in}
% === 错题本排版优化（不参与题目标识） ===
\setlength{\parindent}{0pt}
\setlength{\parskip}{0.35em plus 0.10em minus 0.08em}
\setlength{\emergencystretch}{3em}
\setlist{itemsep=0.25em, topsep=0.40em, parsep=0pt}
\allowdisplaybreaks[2]
% === 错题本排版优化结束 ===
\begin{document}
\section*{首都师范大学2025年高等代数真题}
\section*{29．首都师范大学2025年高等代数试卷}

\subsection{题目1}
1．（15 分）确定以下多项式是否为有理数域 \(\mathbb{Q}\) 上的不可约多项式，详细说明理由．

\begin{itemize}
  \item[（1）] \(x^{p-1}+x^{p-2}+\cdots+x+1\) ，其中 \(p\) 为素数．
  \item[（2）] \(x^{12}+4\) ．
\end{itemize}

\subsection{题目2}
2．（15 分）给定正整数 \(n\) ，实数 \(\alpha \in\left(1, \frac{\pi}{2}\right)\) ，定义一个 \(n\) 阶实方阵 \(A_{n}=\left(a_{i j}\right)_{1 \leq i, j \leq n}\) 满足

\[
a_{i j}= \begin{cases}2 \cos \alpha, & i=j ; \\ 1, & |i-j|=1 ; \\ 0, & \text { 其他情况. }\end{cases}
\]

即

\[
A_{n}=\left(\begin{array}{ccccc}
2 \cos \alpha & 1 & 0 & \cdots & 0 \\
1 & 2 \cos \alpha & 1 & \cdots & 0 \\
0 & 1 & 2 \cos \alpha & \cdots & 0 \\
\vdots & \vdots & \vdots & & \vdots \\
0 & 0 & 0 & \cdots & \cos 2 \alpha
\end{array}\right) .
\]

证明：\(A_{n}\) 的行列式为 \(\left|A_{n}\right|=\frac{\sin ((n+1) \alpha)}{\sin \alpha}\) ．

\subsection{题目3}
3．（15 分）设 \(A=\left(\begin{array}{ccc}-1 & -2 & 6 \\ -1 & 0 & 3 \\ -1 & -1 & 4\end{array}\right)\) ，求 \(A^{2024}\) ．

\subsection{题目4}
4．（15 分）给定正整数 \(n\) ，令 \(V=M_{n}(\mathbb{R})\) 为实数域 \(\mathbb{R}\) 上的 \(n\) 阶方阵构成的 \(\mathbb{R}-\) 线性空间。对\(A, B \in V\) ，定义 \(\langle A, B\rangle=\operatorname{Tr}\left(A^{\mathrm{T}} B\right)\) ，其中 \(A^{\mathrm{T}}\) 表示矩阵 \(A\) 的转置， \(\operatorname{Tr}\left(A^{\mathrm{T}} B\right)\) 表示矩阵 \(A^{\mathrm{T}} B\) 的迹（即矩阵 \(A^{\mathrm{T}} B\) 的对角线元素之和）。\\
（1）证明：\(\langle\),\(\rangle 为 V\) 上的内积。\\
（2）找出 \(V\) 在内积 \(\langle\),\(\rangle 下的一组标准正交基 v_{1}, v_{2}, \cdots, v_{n^{2}}\) ，即这组向量满足

\[
\left(v_{i}, v_{j}\right)= \begin{cases}1, & i=j \\ 0, & i \neq j\end{cases}
\]

证明你的结论。

\subsection{题目5}
5．（15 分）设 \(A\) 为有理数域 \(\mathbb{Q}\) 上的三阶方阵，\(A^{7}=E\) 为三阶单位矩阵。证明：\(A=E\) 。

\subsection{题目6}
6．（15 分）设 \(n\) 为正整数，\(V\) 为实数域 \(\mathbb{R}\) 上的 \(n\) 维线性空间．设 \(\mathcal{A}: V \rightarrow V\) 为 \(V\) 上的线性变换，\(f(\lambda)\) 为 \(\mathcal{A}\) 的特征多项式。设 \(f(\lambda)=g(\lambda) h(\lambda)\) ，其中 \(g(\lambda), h(\lambda)\) 为互素的多项式。令\(V_{g}=\{v \in V \mid g(\mathcal{A}) v=0\}, V_{h}=\{v \in V \mid h(\mathcal{A}) v=0\}\) ，证明：\(V=V_{g} \oplus V_{h}\) ．

\subsection{题目7}
7．（15 分）设 \(n\) 为正整数，\(V\) 为实数域 \(\mathbb{R}\) 上的一个 \(n\) 维线性空间， \(\mathcal{A}: V \rightarrow V\) 为 \(V\) 上的一个线性变换， \(\mathcal{E}\) 为 \(V\) 上的恒等变换，对正整数 \(k\) ，实数 \(\lambda\) ，令 \(V_{\lambda}^{(k)}=\left\{v \in V \mid(\mathcal{A}-\lambda \mathcal{E})^{k} v=0\right\}\) ，则 \(V_{\lambda}^{(k)}\) 为 \(V\)的子空间。对 \(V\) 的子空间 \(W\) ，记 \(\operatorname{dim} W\) 为 \(W\) 的维数。令 \(V_{\lambda}=\sum_{k \geq 1} V_{\lambda}^{(k)}\) 。

\begin{itemize}
  \item[（1）] 证明：除了有限多个 \(\lambda \in \mathbb{R}\) 以外， \(\operatorname{dim} V_{\lambda}=0\) 。
  \item[（2）] 证明：\(\sum_{\lambda \in \mathbb{R}} \operatorname{dim} V_{\lambda}=n\) 当且仅当 \(\mathcal{A}\) 的所有特征根都是实数。
  \item[（3）] 证明：存在 \(V\) 的一组基使得 \(\mathcal{A}\) 在这组基下的矩阵为对角矩阵当且仅当 \(\sum_{\lambda \in \mathbb{R}} \operatorname{dim} V_{\lambda}^{(1)}=n\) 。
\end{itemize}

\subsection{题目8}
8．（15 分）在欧氏空间 \(\mathbb{R}^{4}\) 中，设 \(V\) 是如下方程组的解空间：

\[
\left\{\begin{array}{l}
2 x_{1}-x_{2}+3 x_{3}-x_{4}=0 \\
3 x_{1}-2 x_{2}-2 x_{4}=0 \\
3 x_{1}-x_{2}+9 x_{3}-x_{4}=0
\end{array}\right.
\]

求 \(V\) 的一组基以及 \(V\) 在 \(\mathbb{R}^{4}\) 的正交补 \(V^{\perp}\) 的一组基，并分别对它们作标准正交化。

\subsection{题目9}
9．（15 分）假设 \(A \in M_{3}(\mathbb{C})\) 是一个三阶复方阵，对任意正整数 \(n\) ，令 \(t_{n}=\operatorname{Tr}\left(A^{n}\right)\) 为矩阵 \(A^{n}\) 的迹（即矩阵 \(A^{n}\) 的对角线上元素之和）。证明：\(|A|=\frac{1}{3} t_{3}-\frac{1}{2} t_{1} t_{2}+\frac{1}{6} t_{1}^{3}\) 。

\subsection{题目10}
10．（15 分）求正交矩阵 \(S\) ，使得 \(S^{\mathrm{T}} A S\) 为对角矩阵，其中

\[
A=\left(\begin{array}{ccc}
23 & -8 & -8 \\
-8 & 11 & -16 \\
-8 & -16 & 11
\end{array}\right)
\]

\(S^{\mathrm{T}}\) 表示矩阵 \(S\) 的转置．
\end{document}
